% !TeX root = ../教材题库_全部条目_转写工作稿.tex
% Source: XJ-B1-C05-EX-P183-Y0458
% 第5章 / 5.3 三角函数的图象与性质 / 教材 p.183 / PDF 第190页
% OCR未校对：公式与文字均可能有误；未完成数学TeX转写；禁止直接发给学生
% 原题及相邻原解截图（核对用）：../题目原图/XJ-B1-C05-EX-P183-Y0458-01.webp
% 本文件只在转写完成后作为题干使用；OCR 条目尚不是可靠题干。
\begin{bankitem}{5.3 例题 例题4}{第5章 / 5.3 三角函数的图象与性质 / 教材 p.183 / PDF 第190页}{XJ-B1-C05-EX-P183-Y0458}{OCR未校对：公式与文字均可能有误；未完成数学TeX转写；禁止直接发给学生}
\begin{ocrdraft}
[教材p.183 / PDF第190页]
\par
0互求函数！y一（2，十斧）的单调递增区间．
\par
解令之一2工十一，则y==srn2x十一=srn乙
\par
的单调递增区间是〔一
\par
因为函数3，=sin之
\par
一十2差丌，
\par
你能仿例4求出函数
\par
3一。、（2，+刊与冖
\par
一十2差丌，差e厶于是由
\par
2
\par
（一2，+）的单调递
\par
一十2差丌2工十一一十2差丌，
\par
2
\par
增区间吗？
\par
得一一丌十差丌主十差丌，k\ensuremath{\text{[OCR:EUR]}}-Z.
\par
因此，函数！y一“，（2，十斧）的单调递增区间是
\par
一丌十差丌，一十差丌，k\ensuremath{\text{[OCR:EUR]}}Z.
\end{ocrdraft}
\end{bankitem}
